% Propietario preservado y actualizado por el cierre sucesor de sus obligaciones internas.
% Integración por delta: arquitectura cosmológica con frontera, rebote y
% componentes. No altera los teoremas Cartan--Holst ni el selector de G.

\section{Condiciones de frontera, regularidad del rebote y descomposición de las componentes cosmológicas}
\label{chap:cosmologia-frontera-rebote}
\index{cosmología!realización por hojas}
\index{rebote cosmológico}\index{componente cosmológica descendiente}
\index{cobordismo}\index{información!conservación cosmológica}

Este capítulo distingue dos dominios. La primera sección demuestra las tres
foliaciones de de Sitter. Las secciones posteriores definen las tuplas de
datos necesarias para una realización cosmológica: selector HMT, carta de
frontera, dinámica de rebote y componente descendiente. Cada conclusión se
restringe al subdominio en el que se verifican las ecuaciones de evolución y
las condiciones de empalme declaradas.

\subsection{Correspondencia entre los regímenes elíptico, parabólico e hiperbólico del TRIT y las cartas AdS–euclídea–dS}
\label{sec:atlas-ads-euclideo-ds-hmt}

Las tres álgebras cuadráticas del TRIT reciben una realización geométrica
coordinada. Para escalas \(\ell,H^{-1}>0\), fijamos las cartas
\[
 \mathcal M_-=(\mathbb D, g_{\rm AdS}),
 \qquad
 \mathcal M_0=(\mathscr H_0^{(L)},g_{\rm E}),
 \qquad
 \mathcal M_+=(\mathrm{dS}_4,g_{\rm dS}),
\]
donde la sección espacial hiperbólica se representa en el disco de Poincaré
por
\[
 g_{\rm AdS}^{\rm sp}
 =\frac{4\ell^2\,|dz|^2}{(1-|z|^2)^2},
 \qquad |z|<1,
\]
la L gnomónica porta la carta euclídea \(g_{\rm E}\), y la carta positiva
adopta una de las foliaciones de de Sitter construidas en la sección
siguiente. El functor geométrico conserva la etiqueta de hoja y los mapas de
cambio:
\[
 \mathfrak R_{\rm geo}:\mathcal X_{\rm TPK}^{\rm enr}
 \longrightarrow
 \mathcal M_-\cup\mathcal M_0\cup\mathcal M_+.
\]

La completación EMRH distingue interior, corona y borde, mientras la doble
proyección conserva una misma sección antes de sus evaluaciones. Sobre la
carta negativa, la restricción al borde y la representación de incidencia
forman el cuadrado
\[
\begin{CD}
 \mathcal X_{\rm TPK}^{\rm enr} @>{\mathfrak R_{\rm AdS}}>>
 \mathcal M_-\\
 @V{\mathcal P_{\rm exc}}VV @VV{\operatorname{Res}_{\partial}}V\\
 \mathscr E_{\partial} @>>{\mathfrak R_{\rm conf}}>
 \operatorname{Obs}(\partial\mathcal M_-),
\end{CD}
\]
con \(\mathfrak R_{\rm conf}\circ\mathcal P_{\rm exc}
=\operatorname{Res}_{\partial}\circ\mathfrak R_{\rm AdS}\). Esta identidad
realiza en HMT la correspondencia holográfica interior--borde: el estado del
interior y los observables conformes de frontera son dos representaciones de
la misma genealogía.

La carta euclídea centraliza el descenso local de los lectores inercial y
gravitatorio,
\[
 \mathcal G_L/\mathcal I_L=1,
\]
y constituye la residencia del principio de equivalencia. El estado anterior
a esa proyección conserva la Ambivalencia areal--volumétrica,
\[
 \mathcal L_{\rm ar}/\mathcal L_{\rm vol}
 =\pi_{\rm HMT}/\varphi_{\rm HMT}^{2}.
\]
La carta positiva completa el atlas con \(\ddot a/a=H^2>0\). Por tanto, la
misma terna geométrica contiene la correspondencia holográfica en curvatura
negativa, la equivalencia local en curvatura nula y la aceleración cosmológica
en curvatura positiva.

\subsection{\(3\) foliaciones exactas de de Sitter}
\label{sec:tres-foliaciones-de-sitter}

\begin{proposition}[Foliaciones cerrada, plana y abierta]
\label{prop:tres-foliaciones-de-sitter}
Para \(H>0\), las tres foliaciones FLRW de de Sitter de radio \(H^{-1}\)
se expresan, con \(a_\ast>0\) en la carta plana, mediante
\[
\begin{array}{ccc}
\toprule
k&a_k(t)&\text{dominio}\\
\midrule
+1&H^{-1}\cosh(Ht)&t\in\mathbb R\\
0&a_\ast e^{Ht}&t\in\mathbb R\quad\text{(parche plano)}\\
-1&H^{-1}\sinh(Ht)&t>0\quad\text{(parche abierto)}
\\\bottomrule
\end{array}
\]
y satisfacen
\[
\boxed{
\left(\frac{\dot a_k}{a_k}\right)^2+\frac{k}{a_k^2}=H^2,
\qquad
\frac{\ddot a_k}{a_k}=H^2,}
\]
\[
\boxed{
R^{(4)}=12H^2,\qquad
R^{(3)}=\frac{6k}{a_k^2}.}
\]
\end{proposition}

\begin{proof}
Para \(k=+1\), la primera ecuación utiliza
\(\tanh^2+\operatorname{sech}^2=1\). Para \(k=-1\), utiliza
\(\coth^2-\operatorname{csch}^2=1\). En la carta plana,
\(\dot a/a=H\). Las tres funciones cumplen
\(\ddot a/a=H^2\), y la fórmula FLRW de curvatura da
\(R^{(4)}=12H^2\). Cada rebanada espacial tridimensional tiene curvatura
seccional constante \(k/a_k^2\); en dimensión tres, la contracción del tensor
de curvatura da por tanto
\(R^{(3)}=3(3-1)k/a_k^2=6k/a_k^2\).
\end{proof}

El signo \(k\) tipa la curvatura de las secciones espaciales, no la
curvatura del espacio-tiempo completo: las tres cartas tienen
\(R^{(4)}>0\). La carta cerrada es global; las cartas plana y abierta son
parches.

\begin{statusbox}
    Para una ruta de retorno cerrada y la elección espacial declarada
    \(D=3\), el cociclo HMT de escala aporta los factores exactos
    \(\varphi^3,\varphi^9,\varphi^{36}\) en las profundidades
    \(9,27,108\), junto con una terna espectral de memoria y una pantalla de
    firma \((3,1)\). La selección cosmológica consiste en construir, sobre
    estado enriquecido de TPK, una medida proyectiva, un factor de escala y una ecuación de
    evolución que entrelacen esos objetos con las foliaciones anteriores. Su
    variable de duración debe derivarse de la recta de acción del
    \cref{chap:escala-accion} y de su carta metrológica, no identificarse con
    el mero número de pasos. \(H_0\), las densidades cosmológicas y el signo
    de aceleración no se emplean como entradas del selector.
\end{statusbox}

Las tres foliaciones constituyen el codominio geométrico exacto de
comparación. El selector construido sobre el estado enriquecido determina la
carta realizada por cada régimen, compone sus dominios y conserva los
observables que distinguen la selección. La identificación con
\(k=+1,0,-1\) se efectúa mediante ese selector y no por la mera coincidencia
de tres signos.

\subsection{Universo con frontera}
\label{sec:universo-con-frontera-hmt}

La completación
\[
 1000=729+271=729+243+27+1
\]
distingue interior, estratos de frontera y sello. En una realización
cosmológica, esta estratificación orienta el tipo de una construcción sobre
un par
\[
 (M,\partial M),
\]
dotado de co-tétrada, conexión de Cartan, materia, memoria y flujo. La
identidad cardinal no produce por sí sola una variedad, una métrica, una
co-tétrada ni condiciones de borde. Es una guía de organización que sólo
adquiere contenido físico mediante un mapa
\[
 \mathfrak R_{\rm cos}:
 \mathcal X_{\rm TPK}^{\rm enr}
 \dashrightarrow
 \bigl(M,\partial M,e,\omega,\Psi,\mathcal B_{\partial M}\bigr).
\]

La frontera forma parte constitutiva del problema variacional. Deben
declararse su carácter espacial, temporal o nulo; los términos de borde y de
esquina; los datos que se fijan; y las condiciones de compatibilidad con las
ecuaciones de movimiento. La memoria HMT proporciona un dominio candidato
para registrar historias y prolongaciones, pero su identificación con datos
geométricos exige el entrelazamiento de holonomías del
\cref{crit:puente-tpk-cartan}.

\subsection{Criterio de regularidad del rebote bajo concentración extrema y dominancia torsional}
\label{sec:rebote-cosmologico-condicionado}

Una concentración no acotada de densidad de energía, acción o curvatura no
puede atribuirse a un mecanismo de almacenamiento energético en el sector
torsional. En la acción mínima, la corriente de espín determina localmente la
contorsión; la curvatura conserva su sector radiativo. Una corrección
repulsiva de alta densidad requiere un modelo de materia con espín, un
acoplamiento, una ecuación de estado y un signo efectivo determinados.

En un modelo de Einstein--Cartan con fluido de espín puede aparecer
\[
 s^2\propto a^{-6}.
\]
Si \(\rho_{\rm m}\) denota densidad de masa, una ecuación efectiva puede
adoptar la forma
\[
 H^2
 =
 \frac{8\pi G}{3}
 \bigl(\rho_{\rm m}-\rho_{\rm spin,m}\bigr)
 +\cdots.
\]
Si se emplea densidad de energía \(\rho_{\rm E}\), la conversión dimensional
obliga a escribir
\[
 H^2
 =
 \frac{8\pi G}{3c^2}
 \bigl(\rho_{\rm E}-\rho_{\rm spin,E}\bigr)
 +\cdots.
\]
No se mezclan ambas convenciones en una misma ecuación.

El candidato a rebote se sitúa en
\[
 \rho_{\rm m}=\rho_{\rm spin,m}
\]
o en la igualdad energética equivalente, siempre que la solución sea regular
y satisfaga las condiciones de materia y frontera declaradas. La acción
Cartan--Holst permite estudiar este mecanismo; HMT debe construir todavía el
mapa que seleccione la corriente, el umbral y el registro de información sin
consultarlos retrospectivamente.

\subsection{Cobordismo y componentes cosmológicas descendientes}
\label{sec:cobordismo-universo-hijo}

La apertura de una componente cosmológica descendiente se formula mediante un cobordismo
\[
 \mathcal W:
 \Sigma_{\rm in}
 \longrightarrow
 \Sigma_{\rm prog}\sqcup\Sigma_{\rm desc}.
\]
Antes de afirmar conservación de información debe fijarse el objeto
conservado. Si \(\mathscr I\) es una carga aditiva definida por un cociclo,
puede ensayarse
\[
 \mathscr I(\Sigma_{\rm in})
 =
 \mathscr I(\Sigma_{\rm prog})
 +\mathscr I(\Sigma_{\rm desc})
 +\mathscr F_{\partial\mathcal W}.
\]
Si \(\mathscr I\) representa información cuántica, la estructura pertinente
es una isometría o una evolución unitaria y el balance no se reduce, en
general, a una suma escalar. Han de calcularse las entropías de los
subsistemas, sus correlaciones y la información mutua.

Una realización debe especificar:
\begin{enumerate}
 \item las condiciones de unión entre las secciones inicial, progenitora y
 descendiente;
 \item los términos de borde y esquina de la acción;
 \item la causalidad, la hiperbolicidad y el problema de valores iniciales;
 \item el control de singularidades, curvas temporales cerradas y condiciones
 de energía;
 \item la carga, isometría o ley de conservación concreta que representa la
 información.
\end{enumerate}
Las identidades de Bianchi y el balance de Noether son condiciones locales
necesarias; no demuestran por sí solos el cambio de topología ni la
conservación global de información.

\subsection{Componentes paralelas y modo intercomponente}
\label{sec:componentes-cosmologicas-paralelas}

Una frontera desconectada y una colección de sectores poseen composiciones
distintas. En una realización de tipo topológico,
\[
 \mathcal H_{\Sigma_{\rm parent}\sqcup\Sigma_{\rm child}}
 \simeq
 \mathcal H_{\rm parent}\otimes\mathcal H_{\rm child}
\]
representa la composición de componentes de una frontera. En cambio,
\[
 \mathcal H_{\rm cos}
 =\bigoplus_a\mathcal H_a
\]
representa alternativas o sectores de superselección. El producto tensorial
y la suma directa no se sustituyen mutuamente.

Sobre la suma directa, un operador global candidato tiene forma de bloques:
\[
 \mathbb H_{\rm cos}
 =
 \begin{pmatrix}
 H_1&K_{12}&\cdots\\
 K_{21}&H_2&\cdots\\
 \vdots&\vdots&\ddots
 \end{pmatrix}.
\]
Para definir un operador físico deben fijarse un dominio denso común o
hipótesis de acotación relativa, un reloj común o relacional y la condición
\[
 K_{ba}=K_{ab}^{\dagger}.
\]
Los bloques no diagonales han de derivarse de la frontera, del cobordismo o
de otro acoplamiento declarado. Si son no nulos, las componentes forman
sectores acoplados de una teoría global; no son universos independientes en
sentido fuerte.

Un vector propio con soporte en varios sumandos constituye un modo
intercomponente. Denominar gravitón a su cuanto exige, además, que el sector
físico reducido posea helicidades \(\pm2\), dos polarizaciones transversas,
energía positiva, propagación lumínica, ausencia de fantasmas y taquiones,
identidades de Ward y acoplamiento universal. El soporte global por sí solo
no satisface estas condiciones.

\subsection{Certificación del rebote torsional y obstrucciones exactas}
\label{sec:condiciones-prueba-cosmologia-hmt}

La realización queda certificada por las construcciones siguientes:
\begin{enumerate}
 \item soluciones regulares de rebote con materia y signo derivados;
 \item un problema de frontera bien planteado;
 \item un cobordismo admisible con condiciones de unión;
 \item un transporte de información definido como carga, isometría o
 evolución unitaria;
 \item un operador global autoadjunto con acoplamientos no diagonales
 seleccionados antes de la comparación;
 \item predicciones que distingan la realización de una cosmología de una
 sola componente.
\end{enumerate}

El certificado exige la solución regular, la compatibilidad de las
condiciones de unión, la conservación de la carga o isometría escogida, la
independencia respecto de cartas arbitrarias y la preservación de causalidad
e hiperbolicidad. La energía global se emplea únicamente cuando existe una
simetría temporal o una definición cuasilocal explícita.
